\documentclass[10pt,a4paper]{amsart}
\usepackage[margin=19mm]{geometry}
\usepackage{amsmath,amssymb,mathtools}
\usepackage[scr=rsfs,cal=euler]{mathalfa}
\usepackage{xfrac}
\usepackage[unicode,hidelinks,bookmarks=false]{hyperref}
\newcommand{\ie}{i.e.\ }
\newcommand{\cf}{cf.\ }
\newcommand{\resp}{respectively\ }
\newcommand{\Subsection}{\subsection}
\providecommand{\nobreakdash}{\nobreak\hbox{-}}
\setlength{\emergencystretch}{2em}
\multlinegap0pt
\usepackage{mathtools}
\usepackage{amssymb}
\usepackage{float}
\usepackage{tensor}
\usepackage{tikz}
\usepackage[normalem]{ulem}
\usepackage{cancel}
\newtheorem{lemma}{Lemma}
\newtheorem{corollary}{Corollary}
\title{optimal control subjected to the $1D$ Keller-Segel problem}
\author{F. Guillen-Gonzalez, F. Palmero-Ramos, M.A. Rodriguez-Bellido, G. Tierra}
\date{Source: arXiv:2603.18889v1 (CC BY 4.0)}
\begin{document}
\maketitle

\noindent\emph{Source and licence.} optimal control subjected to the $1D$ Keller-Segel problem. Version arXiv:2603.18889v1 (CC BY 4.0), distributed by arXiv under the Creative Commons Attribution 4.0 International licence, http://creativecommons.org/licenses/by/4.0/, as recorded in the arXiv licence metadata for this exact version and shown on the abstract page https://arxiv.org/abs/2603.18889v1. Full licence notice: \texttt{licenses/arxiv-2603.18889-CC-BY-4.0.txt}.\par\smallskip

\noindent\textit{Editorial scope.} Counted unit: the article's corollary corollary (source0.tex lines 924--925). The article's complete original proof of it is delivered with it (source0.tex lines 927--930, one \texttt{proof} environment of 4 lines). No mathematics is written, repaired or re-derived: the statement and the proof are the article's own lines, in the article's own notation, and no retained line is substituted or re-flowed.  Retained uncounted as the source-local context the proof needs: lines 893--895.  Nothing was omitted from any retained span, so the package carries no omission record.  No retained line contains a citation, so the wrapper carries no bibliography and no bibliography entry.  No retained line contains a figure environment, an image-inclusion command or drawing code, so no image is delivered and the package declares no asset dependency.  Editorial formatting only, in addition: an AMS \texttt{amsart} preamble that restates the article's own declarations and macros used by the retained lines, namely the environments \texttt{lemma}, \texttt{corollary} on separate counters, together with the standard packages those lines need.  The retained environments print in this document as Lemma 1, Corollary 1 in order of appearance, whereas the article prints the same items with the numbering of its own full document.  The complete native source of the article as distributed in the arXiv e-print for this exact version is bundled unchanged as \texttt{sources/196733/source0.tex}.
\begin{lemma}  \label{sol-pos}
 For any $b\in \mathbb{R}^J_+$, there exists $u\in \mathbb{R}^J_+$ such that $A\, u+Ch(v^n)\, u =b$.
 \end{lemma} 

 \begin{corollary}  For any $b\in \mathbb{R}^J$, there exists $u\in \mathbb{R}^J$ such that $A\, u+Ch(v^n)\, u =b$. Consequently, the matrix $A+ Ch(v^n) $ is regular.
  \end{corollary}

  \begin{proof} {\it Step1.} For any $c\in \mathbb{R}^J_-$, there exists $w\in \mathbb{R}^J_-$ such that $A\, w+Ch(v^n)\, w =c$. Indeed, taking $w=-u$ with $u\ge 0$, by using Lemma \ref{sol-pos} one can deduce that any solution of $A\, u+Ch(v^n)\, u =b=-c\ge 0$ satisfy $u=-w\geq0$, hence $w\leq0$.
  
 {\it  Step 2.} For any $b\in \mathbb{R}^J$, we can write $b=b_+ + b_-$, therefore $z=u+w$ satisfies $A\, z+Ch(v^n)\, z =b$, where $u\ge 0$ satisfies $A\, u+Ch(v^n)\, u =b_+$ and  $w\le 0$ satisfies $A\, w+Ch(v^n)\, w =b_-$
  \end{proof} 

\end{document}
